\documentclass[10pt,a4paper]{amsart}
\usepackage[margin=19mm]{geometry}
\usepackage{amsmath,amssymb,mathtools}
\usepackage[scr=rsfs,cal=euler]{mathalfa}
\usepackage{xfrac}
\usepackage[unicode,hidelinks,bookmarks=false]{hyperref}
\newcommand{\ie}{i.e.\ }
\newcommand{\cf}{cf.\ }
\newcommand{\resp}{respectively\ }
\newcommand{\Subsection}{\subsection}
\providecommand{\nobreakdash}{\nobreak\hbox{-}}
\setlength{\emergencystretch}{2em}
\multlinegap0pt
\usepackage{amssymb}
\newtheorem{proposition}{Proposition}
\def\geq{\geqslant}
\title{Tautological Modular Forms of Level Two\\ and Degree Two}
\author{Fabien Cl\'ery, Gerard van der Geer}
\date{Source: arXiv:2605.13300v1 (CC BY 4.0)}
\begin{document}
\maketitle

\noindent\emph{Source and licence.} Tautological Modular Forms of Level Two\\ and Degree Two. Version arXiv:2605.13300v1 (CC BY 4.0), distributed by arXiv under the Creative Commons Attribution 4.0 International licence, http://creativecommons.org/licenses/by/4.0/, as recorded in the arXiv licence metadata for this exact version and shown on the abstract page https://arxiv.org/abs/2605.13300v1. Full licence notice: \texttt{licenses/arxiv-2605.13300-CC-BY-4.0.txt}.\par\smallskip

\noindent\textit{Editorial scope.} Counted unit: the article's proposition proposition (source0.tex lines 1081--1083). The article's complete original proof of it is delivered with it (source0.tex lines 1084--1102, one \texttt{proof} environment of 19 lines). No mathematics is written, repaired or re-derived: the statement and the proof are the article's own lines, in the article's own notation, and no retained line is substituted or re-flowed.  No further source-local context is needed: every cross-reference of every retained line resolves inside the two delivered spans.  Nothing was omitted from any retained span, so the package carries no omission record.  No retained line contains a citation, so the wrapper carries no bibliography and no bibliography entry.  No retained line contains a figure environment, an image-inclusion command or drawing code, so no image is delivered and the package declares no asset dependency.  Editorial formatting only, in addition: an AMS \texttt{amsart} preamble that restates the article's own declarations and macros used by the retained lines, namely the environments \texttt{proposition} on separate counters, together with the standard packages those lines need.  The retained environments print in this document as Proposition 1 in order of appearance, whereas the article prints the same items with the numbering of its own full document.  The complete native source of the article as distributed in the arXiv e-print for this exact version is bundled unchanged as \texttt{sources/200492/source0.tex}.
\begin{proposition}
We have $v_{\pi}={\rm ord}_{H_{\pi}}$.
\end{proposition}

\begin{proof}
Using the values $v_{\pi}(\vartheta_{\pi'})$,
the fact that on the Siegel upper half space  $\mathfrak{H}_2$ (or an \'etale cover of
$\mathcal{A}_2[2]$) a local equation of $H_{\pi}$ is given by $\vartheta_{\pi}$ and
the fact that the function field of $\mathcal{A}_2[2]$ is generated by powers of
$\vartheta_{\pi'}/\vartheta_{\pi^{\prime\prime}}$, we see that
$v_{\pi}\geq {\rm ord}_{H_{\pi}}$. Indeed, we can write
a function $F$ locally as $\vartheta^c_{\pi}\, f$ with $f$ expressed in the $\vartheta_{\pi'}$ with $\pi'\neq \pi$.
But if we replace $F$ by the product $F'$ of $\sigma(F)$
for all $\sigma \in \mathfrak{S}_6$ we get
$$
v_{\pi}(F')=\sum_{\sigma} v_{\pi}(\sigma(F)) \geq 
\sum_{\sigma} {\rm ord}_{H_{\pi}}(\sigma(F))=
{\rm ord}_{H_{\pi}}(F')=v_{\pi}(F')\, ,
$$
where the equality $v_{\pi}(F')={\rm ord}_{H_{\pi}}(F')$ follows because these valuations agree in level $1$.
Thus we see that $v_{\pi}(F)={\rm ord}_{H_{\pi}}(F)$.

\end{proof}

\end{document}
